\section{Recovering fitting accuracy without removing the physical defense}
\label{sec:protected-fitting}

The completed experiment leaves a specific approximation question. Independent
model covariances prevented collapse, and the physical mixture supplied a
valid proposal at every tested update, but four-dimensional filtering error
remained above the earlier repaired-guide result. The independent covariance
may spread polynomial resolution over a region larger than the observation
requires. Alternatively, the number of rows, polynomial degree, pair ranks,
or finite optimization budget may be insufficient. The following construction
tests those explanations while retaining a quantitative covariance floor and
the exact physical importance correction. It is an extension of the pair-TT
method developed in this note, not a claim of source-faithful TT-cross.

\subsection{Observation adaptation with a covariance floor}

Let $x,z,m_t\in\mathbb R^d$ and $A,Q,P_t,R_t,C_t,L_t\in\mathbb R^{d\times d}$.
Let $P_t\succeq0$ denote the finite, checked covariance returned by the guide
repair, including its predictive fallback when quadrature remains unresolved.
Let $m_t$ be the corresponding mean. As in
\eqref{eq:robust-chart}, assume $P_0^{\rm prior}\succ0$, $Q\succ0$, and
finite linear dynamics $A$, and define
\begin{equation}
 R_0=P_0^{\rm prior},\qquad R_t=AR_{t-1}A^\top+Q,\qquad
 C_t(\lambda)=(1-\lambda)P_t+\lambda R_t,
 \quad 0<\lambda\le1.
 \label{eq:a11-blend}
\end{equation}
The chart is $x=m_t+L_tu$, where $L_tL_t^\top=C_t(\lambda)$ is the
Cholesky factorization with positive diagonal. The previous state receives
the corresponding chart at time $t-1$. These are coordinate covariances;
they do not replace the physical covariance in the SGQF joint initializer.

\begin{proposition}[A relative covariance floor]
\label{prop:a11-floor}
Under these assumptions, $C_t(\lambda)\succ0$ and
\begin{equation}
 C_t(\lambda)\succeq\lambda R_t,\qquad
 C_t(\lambda)^{-1}\preceq\lambda^{-1}R_t^{-1},\qquad
 \det C_t(\lambda)\ge\lambda^d\det R_t.
 \label{eq:a11-floor}
\end{equation}
If $s_1,\ldots,s_d$ are the eigenvalues of
$R_t^{-1/2}P_tR_t^{-1/2}$, the eigenvalues of
$R_t^{-1/2}C_tR_t^{-1/2}$ are exactly
$\lambda+(1-\lambda)s_i\ge\lambda$. If also $P_t\preceq M R_t$,
the condition number of this whitened covariance is at most
$[\lambda+(1-\lambda)M]/\lambda$.
\end{proposition}
\begin{proof}
The recursion gives $R_t\succ0$ by induction, since for nonzero $v$ the
term $v^\top Qv$ is strictly positive. Subtracting $\lambda R_t$ from
$C_t$ leaves $(1-\lambda)P_t\succeq0$. Whiten by $R_t^{-1/2}$ to obtain
$\lambda I+(1-\lambda)R_t^{-1/2}P_tR_t^{-1/2}$. Diagonalizing its
symmetric second term proves the eigenvalue formula. Inverting those
positive eigenvalues and undoing the congruence proves the inverse bound.
Their product, multiplied by $\det R_t$, proves the determinant bound.
Under the additional upper bound, $0\le s_i\le M$, which gives the
stated ratio of largest to smallest eigenvalue.
\end{proof}

Thus $\lambda=.25$ retains at least half the independent-model standard
deviation in every whitened direction. For example, if $P_t=.01R_t$,
then $C_t(.25)=.2575R_t$: the chart follows some observation contraction
without inheriting the near-collapse. The blend need not shrink every
direction relative to $R_t$ because $P_t\preceq R_t$ is not assumed.
The bound is relative to the physical model scale; it is not an absolute
floating-point guarantee when that scale itself is ill-conditioned. Checked
factorization and finite-value diagnostics therefore remain necessary.

\subsection{What the pair-TT regression is approximating}

Write $z=x_{t-1}$, $x=x_t$, and let $p_{t-1}^{\rm ret}(z)$ be the
normalized density retained by the preceding TT step. With the current
observation $y_t$, the unnormalized fitting target is
\begin{equation}
 \Gamma_t(z,x)=p_{t-1}^{\rm ret}(z)f_t(x\mid z)g_t(y_t\mid x).
 \label{eq:a11-target}
\end{equation}
This uses the retained approximation, not the unknown exact filtering
density. Subsequent particle weights nevertheless use the physical
transition and likelihood, so a poor retained approximation affects proposal
efficiency rather than silently replacing that importance-sampling target.

Let $u=L_t^{-1}(x-m_t)$ and
$v=L_{t-1}^{-1}(z-m_{t-1})$, and write $\varphi_{2d}(u,v)$ for the
standard Gaussian density. For a fixed positive scale $s_t$, define
\begin{equation}
 b_t(u,v)=\left[
 \frac{\Gamma_t(m_{t-1}+L_{t-1}v,m_t+L_tu)
       |\det L_{t-1}|\,|\det L_t|}
      {s_t\varphi_{2d}(u,v)}\right]^{1/2},\qquad
 \mathcal E(h)=\int(h-b_t)^2\varphi_{2d}\,du\,dv.
 \label{eq:a11-amplitude}
\end{equation}
The two determinants matter: omitting either fits a different density.
The pair-TT groups the variables as $(u_1,v_1),\ldots,(u_d,v_d)$.
Its Gaussian initialization is transformed into these same charts before
projection and coefficient compression.

\begin{proposition}[A bounded design correction]
\label{prop:a11-design}
Let $j_t$ be the normalized pullback of the physical SGQF joint and draw
training rows from $r_\delta=\delta\varphi_{2d}+(1-\delta)j_t$,
$0<\delta\le1$. For a fixed $h$ and $s_t$, with $\mathcal E(h)<\infty$,
\begin{equation}
 0<\frac{\varphi_{2d}}{r_\delta}\le\delta^{-1},\qquad
 \mathbb E_{r_\delta}\left[
 \frac{\varphi_{2d}(U,V)}{r_\delta(U,V)}
       (h(U,V)-b_t(U,V))^2\right]=\mathcal E(h).
 \label{eq:a11-design}
\end{equation}
\end{proposition}
\begin{proof}
The Gaussian component gives $r_\delta\ge\delta\varphi_{2d}>0$.
Multiplying the integrand by the sampling density cancels $r_\delta$
and leaves precisely the integral in \eqref{eq:a11-amplitude}.
\end{proof}

The implementation retains $\delta=.2$, hence a design importance ratio
bounded by five. This is a bound on the design correction, not on the target
amplitude or the variance of the regression loss. Increasing the row count
tests whether the existing design resolves the important regions. The
training scale is estimated on the training panel and then frozen for
validation; the fitted $h$ also depends on those rows. Consequently the
proposition does not assert unbiased finite-sample training error for that
data-dependent pair $(h,s_t)$. Independent validation measures its error
conditional on the fitted choice.

\begin{proposition}[Capacity and initialization errors]
\label{prop:a11-capacity}
Let $\mathcal H_{p,r}$ be pair-TT polynomials of coordinate degree at most
$p$ and pair ranks at most $r$. If $p'\ge p$ and $r'\ge r$, then
\begin{equation}
 \inf_{h\in\mathcal H_{p',r'}}\mathcal E(h)
 \le \inf_{h\in\mathcal H_{p,r}}\mathcal E(h).
 \label{eq:a11-capacity}
\end{equation}
For a normalized Gaussian-joint amplitude $b_G$, let $\Pi_p$ be orthogonal
projection onto the full degree-$p$ tensor polynomial space under
$\varphi_{2d}$. For any compressed initializer $h_G$ in that space,
\begin{equation}
 \|b_G-h_G\|_2^2=
 \|b_G-\Pi_pb_G\|_2^2+\|\Pi_pb_G-h_G\|_2^2,
 \label{eq:a11-projection}
\end{equation}
where the norm is in $L^2(\varphi_{2d})$.
\end{proposition}
\begin{proof}
Zero-padding coefficient and rank indices embeds each old core product in
the larger class without changing its polynomial. The infimum over a
superset cannot increase. For the second assertion,
$b_G-\Pi_pb_G$ is orthogonal to every polynomial in the projected space,
including $\Pi_pb_G-h_G$. Expanding the squared norm eliminates the cross
term and yields the identity.
\end{proof}

The two initialization errors distinguish insufficient degree from coefficient
compression. They concern the Gaussian initializer, not the nonlinear target
$b_t$. Nor does nested capacity guarantee improvement after finitely many
alternating sweeps with an L1 core penalty: the optimization is nonconvex,
and the penalty depends on the chosen core representation. Rank is inert
when $d=1$, because there is only one pair core. These facts motivate separate
row, sweep, rank and degree ablations, with L1 reselected on calibration data.

\subsection{A normalized conditional and its physical correction}

Let $h_t$ be the finite fitted polynomial and $\tau_t>0$ its internal
defensive mass. Define $a_t(v)=\int h_t(u,v)^2\varphi_d(u)\,du$.
\begin{proposition}[Conditional normalization]
\label{prop:a11-conditional}
For each finite $z$ and its corresponding $v$, the density
\begin{equation}
 q_t^{TT}(x\mid z)=
 \frac{h_t(u,v)^2+\tau_t}{a_t(v)+\tau_t}
 \frac{\varphi_d(u)}{|\det L_t|}
 \label{eq:a11-conditional}
\end{equation}
is strictly positive and integrates to one in physical coordinates.
\end{proposition}
\begin{proof}
A finite polynomial is square-integrable under the Gaussian measure, so
$a_t(v)$ is finite and nonnegative. The positive $\tau_t$ makes numerator
and denominator positive. Substitute $dx=|\det L_t|du$; integration of the
numerator gives $a_t(v)+\tau_t$, which cancels the denominator.
\end{proof}

The same normalization must be used by the inverse-CDF sampler and the
reported conditional density. Changing coordinates without changing one of
these calculations would break the importance correction. We now retain
the physical mixture from Section~17:
\begin{equation}
 q_{t,\epsilon}(x\mid z)=(1-\epsilon)q_t^{TT}(x\mid z)
                  +\epsilon f_t(x\mid z),\qquad 0<\epsilon<1.
 \label{eq:a11-mixture}
\end{equation}
At the initial time, replace $f_t(\cdot\mid z)$ by the initial prior.
Choose the component with its stated probability, sample from it, and
evaluate the \emph{whole} mixture at the resulting draw.

\begin{proposition}[Exact correction and a finite second moment]
\label{prop:a11-weights}
Fix an ancestor $z$, let $G(x)=f_t(x\mid z)g_t(y_t\mid x)$, and assume
$0<Z(z)=\int G(x)dx<\infty$ and
$\int f_t(x\mid z)g_t(y_t\mid x)^2dx<\infty$. Then
\begin{equation}
 w(x,z)=\frac{G(x)}{q_{t,\epsilon}(x\mid z)},\quad
 \mathbb E_q w=Z(z),\quad
 0\le w(x,z)\le\frac{g_t(y_t\mid x)}{\epsilon},\quad
 \mathbb E_qw^2\le\frac1\epsilon\int f_tg_t^2\,dx<\infty.
 \label{eq:a11-weights}
\end{equation}
For conditionally independent identically distributed draws at this fixed
ancestor, the importance ESS satisfies
\begin{equation}
 \frac{\ESS_N}{N}
 =\frac{(N^{-1}\sum_iw_i)^2}{N^{-1}\sum_iw_i^2}
 \longrightarrow \frac{Z(z)^2}{\mathbb E_qw^2}
 =\frac1{1+\chi^2(\pi_z\Vert q)},\qquad
 \pi_z=G/Z(z),
 \label{eq:a11-ess}
\end{equation}
almost surely.
\end{proposition}
\begin{proof}
The lower bound $q\ge\epsilon f_t$ gives the pointwise inequality wherever
$f_t>0$ and the second-moment bound by
$G^2/q\le f_tg_t^2/\epsilon$. Direct cancellation in $\int qw$ proves the
mean identity. The strong law applies to both $w$ and $w^2$ under the stated
integrability conditions. Finally,
$\chi^2(\pi_z\Vert q)=\int\pi_z^2/q-1
=\mathbb E_qw^2/Z(z)^2-1$.
\end{proof}

The integrability assumption holds for the present Gaussian transition and
independent stochastic-volatility observations. For
$g(y\mid x)=\prod_j(2\pi\beta^2 e^{x_j})^{-1/2}
\exp[-y_j^2e^{-x_j}/(2\beta^2)]$, $\beta>0$, its square is bounded above by
$(2\pi\beta^2)^{-d}\exp(-\sum_jx_j)$, which has finite expectation under a
finite Gaussian transition. This proves finiteness, not a useful uniform
upper bound over arbitrary ancestors or observations. In the actual filter,
ancestors and incoming weights differ; the implementation multiplies $w$
by the previous particle weight and reports ESS before resampling.
Equation~\eqref{eq:a11-ess} is therefore a fixed-ancestor explanation of
weight dispersion, not a theorem equating the full filter's measured ESS
with one conditional chi-square divergence. The earlier mean ESS of roughly
$304/512$ was an empirical result and is not an ESS floor.

\begin{proposition}[Mixture strength has a convex second-moment objective]
\label{prop:a11-epsilon}
For fixed $z$, $h_t$, charts and $\tau_t$, under the preceding integrability
assumption, put $J(\epsilon)=\int G(x)^2/q_{t,\epsilon}(x\mid z)dx$.
On any closed interval $[a,b]\subset(0,1)$,
\begin{equation}
 J'(\epsilon)=-\int\frac{G^2(f_t-q_t^{TT})}{q_{t,\epsilon}^2}dx,
 \qquad
 J''(\epsilon)=2\int\frac{G^2(f_t-q_t^{TT})^2}{q_{t,\epsilon}^3}dx
 \ge0.
 \label{eq:a11-epsilon}
\end{equation}
\end{proposition}
\begin{proof}
The denominator is affine in $\epsilon$, so differentiating its reciprocal
gives the displayed integrands. Uniformly on $[a,b]$,
$f_t/q_{t,\epsilon}\le1/a$ and
$q_t^{TT}/q_{t,\epsilon}\le1/(1-b)$. Therefore the first two differentiated
integrands are bounded in absolute value by constant multiples of
$G^2/q_{t,\epsilon}\le f_tg_t^2/a$, an integrable function.
Dominated differentiation justifies the calculation. Nonnegativity proves
convexity.
\end{proof}

Convexity does not imply that increasing $\epsilon$ reduces $J$; its first
derivative can have either sign. Nor does it make finite-particle filtering
MSE convex in $\epsilon$, since propagation and resampling intervene.
The experiment consequently evaluates a small offline non-harm curve after
fitting: a stronger physical component is retained only if its calibration
filtering error and low-ESS frequency remain acceptable, followed by a fresh
confirmation. The density floor is the protection; ESS and filtering error
measure its realized cost and benefit.

\subsection{Local analytical derivatives and the experiment they support}

The coordinate blend adds no eigenspace selection or eigenvalue clipping.
For a scalar parameter $\theta$, use a dot for its derivative. Assume the
guide branch, $\lambda$, $\epsilon$, basis and ranks are fixed locally,
and that $P_t,m_t,R_t$ and the fitted coefficients are differentiable with
finite derivatives. These assumptions describe a local calculation; L1
active-set changes and guide-acceptance decisions need not satisfy them.

\begin{proposition}[Analytical chart and mixture derivatives]
\label{prop:a11-derivative}
With the stated local assumptions,
\begin{align}
 \dot R_t&=\dot A R_{t-1}A^\top+A\dot R_{t-1}A^\top
             +AR_{t-1}\dot A^\top+\dot Q,\nonumber\\
 \dot C_t&=(1-\lambda)\dot P_t+\lambda\dot R_t,\qquad
 \dot L_t=L_t\Phi(L_t^{-1}\dot C_t L_t^{-\top}),\nonumber\\
 \dot u&=-L_t^{-1}(\dot m_t+\dot L_tu),
 \label{eq:a11-chart-derivative}
\end{align}
where $\Phi(B)$ keeps the strictly lower-triangular entries of the symmetric
$B$, halves its diagonal, and sets its upper triangle to zero. At fixed
physical $x,z$, if $\dot\tau_t$ is included, the conditional score is
\begin{equation}
 \partial_\theta\log q_t^{TT}=
 \frac{2h_t\dot h_t+\dot\tau_t}{h_t^2+\tau_t}
 -\frac{\dot a_t+\dot\tau_t}{a_t+\tau_t}
 -u^\top\dot u-\operatorname{tr}(L_t^{-1}\dot L_t),
 \label{eq:a11-tt-score}
\end{equation}
where $\dot h_t$ includes coefficient, $u$ and $v$ dependence, and
$\dot a_t$ includes coefficient and $v$ dependence. For the fixed mixture
probability, its score is
\begin{equation}
 \partial_\theta\log q_{t,\epsilon}
 =\omega_{TT}\partial_\theta\log q_t^{TT}
  +\omega_f\partial_\theta\log f_t,
 \quad
 \omega_{TT}=\frac{(1-\epsilon)q_t^{TT}}{q_{t,\epsilon}},\quad
 \omega_f=\frac{\epsilon f_t}{q_{t,\epsilon}},\quad
 \omega_{TT}+\omega_f=1.
 \label{eq:a11-mixture-score}
\end{equation}
\end{proposition}
\begin{proof}
Differentiate the covariance recursion and the fixed-weight blend term by
term. Differentiating $C_t=L_tL_t^\top$ and multiplying by $L_t^{-1}$ and
$L_t^{-\top}$ gives $B=E+E^\top$, where $E=L_t^{-1}\dot L_t$ is lower
triangular. Its unique solution is $E=\Phi(B)$, proving the Cholesky formula.
Differentiating $x=m_t+L_tu$ at fixed $x$ gives the displayed $\dot u$;
the previous chart gives the analogous $\dot v$.
Taking the logarithm of \eqref{eq:a11-conditional}, differentiating each
term, and using $\partial\log\varphi_d(u)=-u^\top\dot u$ and
$\partial\log|\det L_t|=\operatorname{tr}(L_t^{-1}\dot L_t)$ proves
\eqref{eq:a11-tt-score}. The integral defining $a_t$ is a finite sum of
Gaussian polynomial moments; its derivative includes all coefficient and
$v$ dependence. Finally differentiate the sum in
\eqref{eq:a11-mixture} and divide by its positive density.
\end{proof}

These formulas establish smooth local chart and conditional-density
operations on the stated branch. A total derivative of the complete
parameter-dependent filter must also include the fitted-core recursion,
guide moments, initial law, likelihood, incoming weights and the treatment
of discrete decisions. That implementation and its verification are outside
this fitting experiment. Analytical tractability of the new chart alone
does not establish HMC readiness.

The A11 experiment compares the full A10 method with $\lambda=.5$ and
$.25$, additional rows and sweeps, and degree/rank increments. It preserves
the physical mixture and exact weight formula, reselects L1 on calibration
only, and uses untouched sequences for the final comparison. The primary
quantity is reference-based filtering mean error. Fit residuals,
initialization errors and KKT residuals explain mechanisms; ESS, maximum
weight and resampling describe importance-sampling behavior. A stronger
four-dimensional reference ladder addresses the unresolved precision in
A10. The mathematical bounds justify retaining the protections; only this
downstream comparison can determine whether the proposed fitting changes
are useful.

